\chapter[2006 --- Fire et al.]{2006: Andrew Z. Fire, Craig C. Mello}
\label{ch:2006_Fire}

\section*{Main Contribution}

\textit{Discovered RNA interference (RNAi), a mechanism by which double-stranded RNA silences genes---both a powerful research tool and a potential therapeutic strategy.}\cite{nobel-2006,lecture-2006}

\section{Official Citation}

for their discovery of RNA interference — gene silencing by double-stranded RNA

\section{Background and Historical Context}

The discovery of RNA interference (RNAi) by Andrew Z. Fire and Craig C. Mello in 1998 represented a fundamental advance in our understanding of how genes are regulated in living organisms. Their work revealed a previously unknown mechanism by which cells could silence the expression of specific genes, a process that would prove to have far-reaching implications for both basic biology and therapeutic development. The 2006 Nobel Prize in Physiology or Medicine recognized this landmark discovery ``for their discovery of RNA interference---gene silencing by double-stranded RNA.''

To understand the significance of Fire and Mello's contribution, it is necessary to appreciate the state of molecular biology in the late 1990s. The central dogma of molecular biology---the flow of information from DNA to RNA to protein---had been well established, and the mechanisms of transcription and translation were understood in considerable detail. However, the regulation of gene expression at the post-transcriptional level remained poorly understood, particularly the mechanisms by which cells could specifically silence or suppress the expression of individual genes.

In 1990, a puzzling phenomenon was observed by Richard Jorgensen and colleagues in their studies of petunia flowers. When they introduced additional copies of the chalcone synthase gene---a gene involved in pigment production---into petunia plants, rather than intensifying the purple color as expected, many of the transgenic plants produced flowers with white or variegated patterns. The extra gene copies appeared to be silencing not only themselves but also the endogenous copy of the same gene. Jorgersen termed this phenomenon ``co-suppression,'' and a similar process was independently discovered in fungi by Romano and Macino, who called it ``quelling.'' These observations were initially puzzling because they did not fit neatly into existing models of gene regulation.

A parallel line of research involving the nematode \textit{Caenorhabditis elegans} was providing additional clues. In 1995, Su Guo and Kenneth Kemphues injected antisense RNA molecules complementary to the \textit{unc-22} gene into \textit{C. elegans} embryos. They observed that the injected worms exhibited a characteristic twitching phenotype, suggesting that the gene had been silenced. However, the mechanism by which antisense RNA produced this silencing effect was not understood. Importantly, control experiments injecting ``scrambled'' sense RNA---which should not have been able to base-pair with any cellular RNA---also produced silencing effects, adding to the confusion.

Andrew Fire, working at the Carnegie Institution of Washington, and Craig Mello, then at the University of Massachusetts Medical School, began collaborating to investigate these puzzling gene-silencing phenomena in \textit{C. elegans}. Their approach was methodical and creative. They tested various combinations of sense and antisense RNA molecules and measured their effects on gene expression. In a series of experiments conducted in 1998, they made the crucial discovery that the silencing activity was not primarily due to either sense or antisense RNA alone, but rather to double-stranded RNA (dsRNA) molecules that combined both strands. When they injected dsRNA corresponding to a specific gene into \textit{C. elegans}, the resulting gene silencing was remarkably potent and specific.

\section{The Discovery/Contribution}

The key experiment that established the discovery of RNA interference was elegantly simple in design but significant in its implications. Fire and Mello prepared RNA molecules of various types targeting the \textit{unc-22} gene in \textit{C. elegans}: single-stranded sense RNA, single-stranded antisense RNA, and double-stranded RNA formed by hybridizing sense and antisense strands together. When injected into the worms, the double-stranded RNA proved to be far more effective at silencing the target gene than either single-stranded form alone. As little as a single molecule of dsRNA per cell was sufficient to produce silencing, suggesting that the mechanism involved some form of catalytic amplification.

Fire and Mello published their landmark paper in \textit{Nature} in February 1998, titled ``Potent and specific genetic interference by double-stranded RNA in \textit{Caenorhabditis elegans}.'' In this paper, they proposed that the introduction of dsRNA into \textit{C. elegans} triggered a sequence-specific degradation of messenger RNA (mRNA) molecules bearing the same sequence, thereby preventing the production of the corresponding protein. They coined the term ``RNA interference'' to describe this phenomenon and recognized its potential as a powerful tool for studying gene function.

The mechanism of RNA interference, elucidated through subsequent research by Fire, Mello, and many other investigators, involves a multi-step pathway. The process begins when a long double-stranded RNA molecule is recognized and cleaved by an enzyme called Dicer, a member of the RNase III family of endonucleases. Dicer cuts the dsRNA into small fragments of approximately 21--23 nucleotides in length, known as small interfering RNAs (siRNAs). These siRNAs consist of two complementary strands with 2-nucleotide 3' overhangs on each strand.

The siRNA duplex is then incorporated into a multi-protein complex called the RNA-induced silencing complex (RISC). Within RISC, one strand of the siRNA---the ``passenger strand''---is degraded, while the other strand---the ``guide strand''---is retained and serves as a template for sequence-specific recognition of target mRNA molecules. When the guide strand finds a complementary sequence in an mRNA molecule, the RISC complex cleaves the mRNA at a specific position within the region of complementarity, leading to its degradation and the prevention of protein production from that transcript.

One of the notable features of RNA interference is its ability to spread systemically throughout the organism. In \textit{C. elegans}, the silencing signal can spread from cell to cell and even from tissue to tissue, allowing a small amount of injected dsRNA to silence gene expression throughout the entire animal. This systemic spread involves the transport of siRNAs through channels called gap junctions and through the extracellular environment. The amplification of the silencing signal involves the activity of RNA-dependent RNA polymerases (RdRPs), which synthesize secondary siRNAs using the target mRNA as a template. This amplification mechanism explains why RNAi in \textit{C. elegans} is so potent and persistent.

The specificity of RNA interference is governed by the sequence complementarity between the siRNA guide strand and the target mRNA. This specificity is precise enough to distinguish between closely related genes and even between allelic variants differing by a single nucleotide. This notable sequence specificity made RNAi an invaluable tool for functional genomics, allowing researchers to systematically knock down individual genes and study their effects on cellular processes and organismal phenotypes.

Fire and Mello's discovery also revealed important connections to previously observed gene-silencing phenomena. Co-suppression in plants and quelling in fungi were subsequently shown to be manifestations of the same fundamental RNAi pathway, involving conserved Dicer and Argonaute proteins. The discovery unified these disparate observations under a single mechanistic framework and revealed that RNA interference was an evolutionarily conserved pathway present in organisms ranging from plants and fungi to insects and mammals.

\section{Impact on Modern Medicine}

The impact of Fire and Mello's discovery on biomedical research and medicine has been transformative. RNA interference rapidly became one of the most widely used tools in molecular biology, enabling researchers to selectively silence genes in virtually any organism and study the resulting phenotypes. This capability revolutionized the field of functional genomics, allowing systematic genome-wide screens to identify genes involved in specific biological processes, disease pathways, and drug responses.

The development of RNAi as a therapeutic strategy has been one of the most exciting areas of clinical translation following Fire and Mello's discovery. The ability to silence disease-causing genes with high specificity offered a fundamentally new approach to treating a wide range of genetic and acquired diseases. The first RNAi therapeutic to receive regulatory approval was patisiran (Onpattro), developed by Alnylam Pharmaceuticals, which was approved by the U.S. Food and Drug Administration (FDA) in 2018 for the treatment of hereditary transthyretin amyloidosis (hATTR), a rare genetic disease caused by the production of misfolded transthyretin protein.

Patisiran works by delivering siRNAs targeting the mRNA encoding transthyretin (TTR) using lipid nanoparticle technology. By silencing the TTR gene in the liver, patisiran reduces the production of both normal and mutant TTR protein, thereby preventing the accumulation of amyloid deposits that cause nerve damage and heart disease. Clinical trials demonstrated that patisiran significantly improved neurological function and quality of life in patients with hATTR amyloidosis.

Following the success of patisiran, additional RNAi therapeutics have been developed and approved for various conditions. Givosiran (Givlaari) was approved in 2019 for the treatment of acute hepatic porphyria, a rare metabolic disorder. Lumasiran (Oxlumo) was approved in 2020 for primary hyperoxaluria type 1, a rare genetic kidney disease. Inclisiran (Leqvio) was approved in 2021 for the treatment of hypercholesterolemia, targeting the PCSK9 gene to lower LDL cholesterol levels. These approvals demonstrated the versatility of RNAi as a therapeutic platform and validated the decades of research that followed Fire and Mello's initial discovery.

The delivery of RNAi therapeutics has been a major technical challenge, as siRNAs are rapidly degraded in the bloodstream and cannot easily cross cell membranes. Advances in chemical modification, lipid nanoparticle formulation, and conjugation strategies---particularly the development of N-acetylgalactosamine (GalNAc) conjugation for liver-targeted delivery---have addressed many of these challenges and enabled the clinical success of RNAi drugs. Ongoing research aims to extend the reach of RNAi therapeutics to other organs and tissues, including the central nervous system, lungs, and muscles.

Beyond direct therapeutic applications, RNAi has become an indispensable tool in drug discovery and development. RNAi-based screens are used to identify and validate drug targets, study drug resistance mechanisms, and develop biomarkers for disease diagnosis and prognosis. The high-throughput nature of RNAi screening allows researchers to systematically evaluate the effects of silencing thousands of genes in cell-based assays, accelerating the identification of novel therapeutic targets.

In basic research, RNA interference has transformed our understanding of gene regulation and cellular processes. Genome-wide RNAi screens in organisms ranging from \textit{C. elegans} to human cells have identified genes involved in virtually every aspect of cell biology, from cell division and metabolism to signaling pathways and disease mechanisms. The discovery of RNAi also led to the recognition of microRNAs (miRNAs) as a major class of endogenous gene regulators, opening up entirely new fields of research in gene regulation and developmental biology.

The Nobel Prize in Physiology or Medicine was awarded to Andrew Z. Fire and Craig C. Mello in 2006 ``for their discovery of RNA interference---gene silencing by double-stranded RNA.'' This recognition highlighted the fundamental importance of their discovery and its impact on our understanding of gene regulation and its potential for therapeutic development.

\section{Legacy}

The legacy of Andrew Fire and Craig Mello's discovery of RNA interference extends far beyond its immediate applications in research and medicine. Their work revealed a fundamental mechanism of gene regulation that is conserved across virtually all eukaryotic organisms, from plants and fungi to insects and humans. This universality underscores the biological importance of RNAi and its deep evolutionary roots.

Fire and Mello's discovery also contributed to a paradigm shift in our understanding of the role of RNA in biology. Rather than serving merely as an intermediate carrier of genetic information between DNA and protein, RNA was revealed to play active regulatory roles in gene expression. This insight catalyzed a broader appreciation of the complexity and diversity of RNA molecules and their functions, contributing to the discovery of other non-coding RNA species, including microRNAs, long non-coding RNAs, and circular RNAs.

The technical innovations inspired by RNAi research have had a lasting impact on biotechnology and pharmaceutical development. The chemical modifications, delivery technologies, and formulation strategies developed for RNAi therapeutics have also enabled the development of other RNA-based medicines, including antisense oligonucleotides and, most notably, messenger RNA (mRNA) therapeutics. The remarkable success of mRNA vaccines against COVID-19, developed by Moderna and BioNTech/Pfizer, drew upon many of the advances in RNA chemistry and delivery that were spurred by RNAi research.

Fire and Mello's work also exemplifies the value of curiosity-driven basic research. Their initial investigations into gene silencing in \textit{C. elegans} were motivated by fundamental questions about how genes are regulated, not by the goal of developing new therapies. Yet the practical applications of their discovery have proven to be immense, illustrating the unpredictable but significant ways in which basic research can lead to transformative medical advances.

Today, RNA interference continues to be an active area of research, with ongoing efforts to improve the efficacy, specificity, and delivery of RNAi-based therapeutics. New applications are being explored in diverse areas, including infectious diseases, cancer, cardiovascular disease, and neurodegenerative disorders. The story of RNAi discovery and its translation into clinical practice stands as one of the great success stories of modern biomedical research, demonstrating the power of fundamental scientific discovery to change the world.


\section{Modern Developments}

Fire and Mello's RNAi discovery has led to modern RNA therapeutics. Understanding of gene silencing has enabled development of siRNA drugs (patisiran for hereditary transthyretin amyloidosis, givosiran for acute hepatic porphyria). RNAi-based pesticides target specific agricultural pests. Understanding of miRNA regulation has led to research on miRNA-based therapies for cancer and cardiovascular disease. RNAi screens identify drug targets and study gene function.


\section{Bibliography}

\begin{thebibliography}{9}
\bibitem{nobel-2006}
Nobel Prize Outreach, ``The Nobel Prize in Physiology or Medicine 2006,''\emph{NobelPrize.org}.\\
\url{https://www.nobelprize.org/prizes/medicine/2006/summary/}

\bibitem{lecture-2006}
Andrew Z. Fire, ``Nobel Lecture,''\emph{NobelPrize.org}. Winner-authored primary source; downloadable from the official Nobel Prize archive.\\
\url{https://www.nobelprize.org/prizes/medicine/2006/fire/lecture/}
\end{thebibliography}
